% Unit 6 (Thermochemistry) guided notes -- 5 blocks, CED sequence.
\documentclass{apchem-notes}

\apcunit{6}
\apcunittitle{Thermochemistry}
\apcdoctitle{Guided Notes}
\apcsubtitle{CED 6.1--6.9 \textbullet\ Zumdahl Ch 6, \S8.8, \S10.8 \textbullet\ 5 blocks}

\begin{document}
\apctitleblock

\begin{apcnote}[The one habit this unit runs on]
Nearly every error in thermochemistry is a \textbf{sign} error, and nearly
every sign error comes from losing track of \emph{whose} point of view you
are taking.

Fix the convention now and use it without exception: signs are always
written \textbf{from the system's point of view}. Heat entering the system
is positive; heat leaving is negative. When a reaction warms the water in a
calorimeter, the water gained heat ($q_{\text{water}} > 0$) and the reaction
lost it ($q_{\text{rxn}} < 0$). Those are the same event described from two
sides, and the exam will ask for both.
\end{apcnote}

% =====================================================================
\blocklesson{Energy, Heat, and Direction}{CED 6.1--6.3 \textbullet\ Zumdahl \S6.1}

\begin{objectives}
  \item Classify processes as endothermic or exothermic and assign signs.
  \item Explain heat transfer and thermal equilibrium at the particle level.
\end{objectives}

\reading{Zumdahl \S6.1}{283--290}

\begin{warmup}[4]
  \item Units of energy: \answerblank[3cm]{joules (J)}
  \item Temperature measures the average
        \answerblank[3.4cm]{kinetic energy} of particles.
  \item Breaking a bond requires energy or releases it?
        \answerblank[2.4cm]{requires}
  \item Which has more energy, a gas or the same substance as a liquid?
        \answerblank[1.8cm]{gas}
\end{warmup}

\segment{INSTRUCTION A}{25}{System, surroundings, and sign}

\notesection{The convention that prevents most errors}{6.1}
\practice{5}

The \term{system} is the part you are studying --- usually the reaction. The
\term{surroundings} are everything else, including the solvent, the
calorimeter, and the thermometer.

\begin{center}
\begin{tabular}{@{}llll@{}}
\hline
\textbf{Process} & \textbf{Heat} & $q_{\text{system}}$ &
  \textbf{Surroundings} \\
\hline
\term{Exothermic} & released by the system & \textbf{negative} & warm up \\
\term{Endothermic} & absorbed by the system & \textbf{positive} & cool down \\
\hline
\end{tabular}
\end{center}

\begin{apctrap}
\textbf{The thermometer measures the surroundings, not the system.} A
reaction that makes the beaker feel hot is \emph{exothermic}, so its
$\Delta H$ is \answerblank[2cm]{negative} --- even though the temperature
you measured went \emph{up}.

This inversion is the most-missed idea in the unit. When a question gives
you a temperature change, your first move should be to write down which
thing changed temperature.
\end{apctrap}

\notesection{Energy diagrams}{6.2}

An energy diagram plots potential energy against reaction progress.

\begin{itemize}[leftmargin=*,itemsep=0.3em]
  \item Products \emph{below} reactants $\Rightarrow$
        \answerblank[2.4cm]{exothermic}, $\Delta H < 0$.
  \item Products \emph{above} reactants $\Rightarrow$
        \answerblank[2.7cm]{endothermic}, $\Delta H > 0$.
  \item $\Delta H$ is the difference between
        \answerblank[4.6cm]{products and reactants} --- and it says nothing
        about how fast the reaction goes.
\end{itemize}

\segment{INSTRUCTION B}{20}{Heat transfer and thermal equilibrium}

\notesection{What actually moves}{6.3}
\practice{6}

Put a hot metal block into cool water. At the particle level:

\begin{itemize}[leftmargin=*,itemsep=0.3em]
  \item Fast-moving metal atoms collide with slower water molecules and
        \answerblank[4.8cm]{transfer kinetic energy}.
  \item Energy flows from \answerblank[2.6cm]{hot to cold} --- always, and
        never the reverse on its own.
  \item Transfer continues until both reach the same
        \answerblank[2.6cm]{temperature}: \term{thermal equilibrium}.
  \item At equilibrium the transfer has not \emph{stopped} --- collisions
        continue, but energy now flows both ways at
        \answerblank[2.4cm]{equal rates}.
\end{itemize}

\begin{apcnote}[Heat is not a substance a body contains]
Say ``heat is transferred,'' not ``the block has a lot of heat.'' Heat is
energy \emph{in transit} because of a temperature difference. What a body
contains is internal energy; heat is what crosses the boundary.

The exam does penalize ``the hot object gives its heat to the cold object''
when the question asks for a particle-level explanation --- it wants
collisions and kinetic-energy transfer.
\end{apcnote}

\segment{APPLICATION}{20}{Signs and directions}

\begin{enumerate}[leftmargin=*,itemsep=0.4em]
  \item Ammonium nitrate dissolves and the beaker becomes noticeably cold.
        Give the sign of $q$ for the dissolving process, the sign for the
        water, and classify it. \answerlines[3]{The dissolving process
        absorbed energy, so $q_{\text{system}} > 0$ --- endothermic. The
        water lost that energy, so $q_{\text{water}} < 0$, which is why it
        cooled. The two have opposite signs and equal magnitude.}
  \item Two blocks of the same metal, one at \SI{80}{\celsius} and one at
        \SI{20}{\celsius}, are placed in contact and left. Describe the
        final state and what is happening in it. \answerlines[3]{Both reach
        the same intermediate temperature --- thermal equilibrium. Energy
        transfer has not ceased: collisions continue at the interface, but
        energy now flows in both directions at equal rates, so no net change
        occurs.}
\end{enumerate}

\begin{exitticket}
A reaction is run in solution and the thermometer reading rises from
\SI{22.0}{\celsius} to \SI{28.4}{\celsius}. State the sign of
$q_{\text{water}}$, the sign of $\Delta H_{\text{rxn}}$, and explain why
they differ.
\keyonly{\begin{apcanswer}$q_{\text{water}}$ is positive --- the water
absorbed energy and warmed. $\Delta H_{\text{rxn}}$ is negative: the
reaction released that energy, so the process is exothermic. They differ
because the two describe the same energy from opposite sides of the system
boundary.\end{apcanswer}}
\end{exitticket}

% =====================================================================
\blocklesson{Heat Capacity and Calorimetry}{CED 6.4 \textbullet\ Zumdahl \S6.2}

\begin{objectives}
  \item Use $q = mc\Delta T$ in both directions.
  \item Determine $\Delta H$ per mole from calorimetry data.
\end{objectives}

\reading{Zumdahl \S6.2}{290--297}

\begin{warmup}[3]
  \item Exothermic means $q_{\text{system}}$ is:
        \answerblank[2cm]{negative}
  \item Heat flows from: \answerblank[2.6cm]{hot to cold}
  \item Specific heat of water: \answerblank[3.4cm]{\SI{4.18}{\joule\per\gram\per\celsius}}
\end{warmup}

\segment{INSTRUCTION A}{25}{The workhorse equation}

\notesection{$q = mc\Delta T$}{6.4}
\practice{5}

\[ q = mc\Delta T \qquad \Delta T = T_{\text{final}} - T_{\text{initial}} \]

\term{Specific heat} $c$ is the energy needed to raise \SI{1}{\gram} by
\SI{1}{\celsius}. Water's is unusually
\answerblank[2cm]{large} at \SI{4.18}{\joule\per\gram\per\celsius}, which is
why it is used as the absorbing medium in calorimetry and why coastal
climates are mild.

\begin{workedexample}[Worked example 1: straight substitution]
How much energy raises \SI{50.0}{\gram} of water from \SI{20.0}{\celsius}
to \SI{80.0}{\celsius}?
\[ q = (50.0)(4.18)(60.0) = 12{,}540~\text{J} =
   \mathbf{12.5}~\text{kJ} \]
Positive, because the water absorbed energy.
\end{workedexample}

\segment{INSTRUCTION B}{20}{Calorimetry: from $q$ to $\Delta H$ per mole}

\notesection{The two-step chain}{6.4}
\practice{5}

A calorimeter measures the temperature change of the \emph{water}. Getting
to $\Delta H$ takes two moves:

\begin{enumerate}[leftmargin=*,itemsep=0.3em]
  \item $q_{\text{water}} = mc\Delta T$, then $q_{\text{rxn}}$ equals
        \answerblank[3cm]{$-q_{\text{water}}$}.
  \item Divide by the moles of the limiting reactant to get
        \answerblank[3cm]{\si{\kilo\joule\per\mole}}.
\end{enumerate}

\begin{workedexample}[Worked example 2: a full calorimetry problem]
A reaction is run in \SI{100.0}{\gram} of water in a coffee-cup
calorimeter. The temperature rises by \SI{5.20}{\celsius}, and
\SI{0.0250}{\mole} of the limiting reactant was used.

\textbf{Step 1 --- heat absorbed by the water:}
\[ q_{\text{water}} = (100.0)(4.18)(5.20) = 2174~\text{J} \]
The water \emph{gained} it, so the reaction \emph{lost} it:
$q_{\text{rxn}} = -2174$~J.

\textbf{Step 2 --- per mole:}
\[ \Delta H = \frac{-2174~\text{J}}{0.0250~\text{mol}}
   = -86{,}900~\si{\joule\per\mole} =
   \mathbf{-86.9}~\si{\kilo\joule\per\mole} \]

Negative, as the temperature rise already told us.
\end{workedexample}

\begin{apctrap}
Two habits worth building here. First, \textbf{the minus sign in step 1 is
not optional} --- omitting it reports an exothermic reaction as
endothermic. Second, \textbf{$\Delta H$ is per mole}; a bare answer in
kilojoules with no ``per mole'' loses the point even when the arithmetic is
right.

Note also that we used the \emph{water's} mass and specific heat, not the
reactants'. The water is the surroundings, and it is what the thermometer
was in.
\end{apctrap}

\segment{APPLICATION}{20}{Thermal equilibrium calculations}

\begin{enumerate}[leftmargin=*,itemsep=0.4em]
  \item A \SI{50.0}{\gram} metal block at \SI{100.0}{\celsius} is dropped
        into \SI{100.0}{\gram} of water at \SI{22.0}{\celsius}. The final
        temperature is \SI{24.6}{\celsius}. Find the specific heat of the
        metal. \workspace[3cm]
        \keyonly{\begin{apcanswer}$q_{\text{water}} = (100.0)(4.18)(2.6) =
        1087$~J, gained. So the metal lost 1087~J:
        $-1087 = (50.0)(c)(24.6 - 100.0) = (50.0)(c)(-75.4)$,
        giving $c = \mathbf{0.288}~\si{\joule\per\gram\per\celsius}$.
        Much smaller than water's --- typical of a
        metal.\end{apcanswer}}
  \item Why does the water's temperature change so little while the metal's
        changes by more than \SI{75}{\celsius}?
        \answerlines[3]{Water has a far larger specific heat, so absorbing
        a given amount of energy raises its temperature much less. There is
        also twice as much water by mass. Both factors work the same way.}
\end{enumerate}

\begin{exitticket}
In a calorimetry problem, whose mass and whose specific heat go into
$q = mc\Delta T$, and why?
\keyonly{\begin{apcanswer}The water's --- the surroundings'. The
calorimeter measures the temperature change of the water, and that is the
only thing whose heat we can compute directly. The reaction's enthalpy is
then inferred as the negative of that value, divided by
moles.\end{apcanswer}}
\end{exitticket}

% =====================================================================
\blocklesson{Energy of Phase Changes}{CED 6.5 \textbullet\ Zumdahl \S10.8}

\begin{objectives}
  \item Calculate the energy of a phase change.
  \item Read a heating curve and choose the right equation for each region.
\end{objectives}

\reading{Zumdahl \S10.8}{523--530}

\begin{warmup}[3]
  \item $q = mc\Delta T$ applies when what is changing?
        \answerblank[2.6cm]{temperature}
  \item $\Delta H_{\text{fus}}$ for water:
        \answerblank[3cm]{\SI{6.02}{\kilo\joule\per\mole}}
  \item $\Delta H_{\text{vap}}$ for water:
        \answerblank[3cm]{\SI{40.7}{\kilo\joule\per\mole}}
\end{warmup}

\segment{INSTRUCTION A}{25}{Temperature does not change during a phase change}

\notesection{Two different equations}{6.5}
\practice{5}

\begin{center}
\fbox{\parbox{0.9\linewidth}{\centering
\textbf{Temperature changing, one phase:} $q = mc\Delta T$ \\[0.3em]
\textbf{Phase changing, temperature constant:}
$q = n\Delta H_{\text{fus}}$ or $n\Delta H_{\text{vap}}$}}
\end{center}

\begin{apctrap}
\textbf{You cannot use $q = mc\Delta T$ during a phase change}, because
$\Delta T$ is \answerblank[1.4cm]{zero} --- it would give $q = 0$, which is
plainly wrong for melting a block of ice.

During melting or boiling, the added energy goes into overcoming
\textbf{intermolecular forces}, not into speeding the particles up. That is
why the temperature holds steady on a heating curve, and it is the
explanation the exam wants.
\end{apctrap}

\begin{workedexample}[Worked example 3: melting and boiling]
\textbf{Melt \SI{36.0}{\gram} of ice at \SI{0}{\celsius}.}
$n = 36.0/18.02 = \SI{2.00}{\mole}$;
$q = (2.00)(6.02) = \mathbf{12.0}~\text{kJ}$.

\textbf{Vaporize \SI{18.0}{\gram} of water at \SI{100}{\celsius}.}
$n = \SI{1.00}{\mole}$; $q = (1.00)(40.7) = \mathbf{40.7}~\text{kJ}$.

Vaporizing one mole costs nearly \textbf{seven times} as much as melting
one mole --- because melting only loosens the lattice, while boiling
separates the molecules entirely.
\end{workedexample}

\segment{INSTRUCTION B}{20}{Reading a heating curve}

\notesection{Five regions, five calculations}{6.5}
\practice{4}

\begin{center}
\begin{tikzpicture}
\begin{axis}[width=10.6cm, height=5cm,
    xlabel={energy added}, ylabel={temperature},
    xtick=\empty, ytick=\empty, xmin=0, xmax=12, ymin=-2, ymax=12,
    axis lines=left]
  \addplot[seriesA] coordinates
    {(0,-1.5) (1,0) (4,0) (5.5,6) (10.5,6) (11.5,9.5)};
  \node[font=\scriptsize,anchor=west] at (axis cs:0.1,1.2) {solid};
  \node[font=\scriptsize,anchor=south] at (axis cs:2.5,0.2) {melting};
  \node[font=\scriptsize,anchor=west] at (axis cs:4.6,3.4) {liquid};
  \node[font=\scriptsize,anchor=south] at (axis cs:7.5,6.2) {boiling};
  \node[font=\scriptsize,anchor=west] at (axis cs:10.6,8.2) {gas};
\end{axis}
\end{tikzpicture}
\end{center}

The \textbf{sloped} regions use \answerblank[2.6cm]{$q = mc\Delta T$}; the
\textbf{flat} regions use \answerblank[3cm]{$q = n\Delta H$}.

The boiling plateau is always much \answerblank[2cm]{longer} than the
melting plateau, for the same reason vaporization costs more energy.

\begin{workedexample}[Worked example 4: ice to steam, all five regions]
Take \SI{25.0}{\gram} of ice from \SI{-10.0}{\celsius} to steam at
\SI{110.0}{\celsius}. ($c_{\text{ice}} = 2.09$,
$c_{\text{water}} = 4.18$, $c_{\text{steam}} = 2.03$
\si{\joule\per\gram\per\celsius}; $n = 25.0/18.02 = \SI{1.39}{\mole}$)

\begin{center}
\begin{tabular}{@{}llr@{}}
\hline
\textbf{Region} & \textbf{Equation} & \textbf{$q$ (kJ)} \\
\hline
warm ice $-10 \to 0$ & $mc\Delta T$ & 0.52 \\
melt at \SI{0}{\celsius} & $n\Delta H_{\text{fus}}$ & 8.35 \\
warm water $0 \to 100$ & $mc\Delta T$ & 10.45 \\
boil at \SI{100}{\celsius} & $n\Delta H_{\text{vap}}$ & \textbf{56.47} \\
warm steam $100 \to 110$ & $mc\Delta T$ & 0.51 \\
\hline
\textbf{Total} & & \textbf{76.30} \\
\hline
\end{tabular}
\end{center}

Boiling alone is \textbf{74\%} of the total. If a heating-curve question
asks which step dominates, that is nearly always the answer.
\end{workedexample}

\segment{APPLICATION}{20}{Phase-change reasoning}

\begin{enumerate}[leftmargin=*,itemsep=0.4em]
  \item Calculate the energy to melt \SI{90.0}{\gram} of ice at
        \SI{0}{\celsius}. \workspace[2cm]
        \keyonly{\begin{apcanswer}$n = 90.0/18.02 = \SI{4.99}{\mole}$;
        $q = (4.99)(6.02) = \mathbf{30.1}~\text{kJ}$\end{apcanswer}}
  \item Explain, in terms of intermolecular forces, why the temperature
        stays constant while ice melts even though energy is being added.
        \answerlines[3]{The added energy is used to overcome the hydrogen
        bonds holding the crystal together, increasing the potential energy
        of the particles. Temperature reflects average \emph{kinetic}
        energy, which is not changing, so the thermometer reading holds
        steady until the last of the solid is gone.}
  \item Why is a steam burn worse than a burn from the same mass of water
        at \SI{100}{\celsius}? \answerlines[3]{Condensing steam releases
        \SI{40.7}{\kilo\joule\per\mole} before its temperature falls at
        all. That energy is delivered to the skin on top of whatever the
        subsequent cooling releases, so far more energy is transferred.}
\end{enumerate}

\begin{exitticket}
Which region of a heating curve requires the most energy for water, and
why?
\keyonly{\begin{apcanswer}Boiling. Vaporization must separate the molecules
completely against their hydrogen bonding, whereas melting only has to
loosen the lattice --- $\SI{40.7}{\kilo\joule\per\mole}$ against
$\SI{6.02}{\kilo\joule\per\mole}$, nearly seven times
larger.\end{apcanswer}}
\end{exitticket}

% =====================================================================
\blocklesson{Enthalpy of Reaction and Hess's Law}{CED 6.6, 6.9 \textbullet\ Zumdahl \S6.2--6.3}

\begin{objectives}
  \item Scale $\Delta H$ with the amount of substance reacting.
  \item Combine reactions using Hess's law.
\end{objectives}

\reading{Zumdahl \S6.2--6.3}{290--301}

\begin{warmup}[3]
  \item During a phase change, $\Delta T$ is:
        \answerblank[1.6cm]{zero}
  \item $\Delta H_{\text{vap}}$ is larger than
        $\Delta H_{\text{fus}}$ because:
        \answerblank[7.5cm]{boiling separates molecules entirely}
  \item Exothermic $\Delta H$ is: \answerblank[2cm]{negative}
\end{warmup}

\segment{INSTRUCTION A}{25}{$\Delta H$ is an extensive quantity}

\notesection{Enthalpy scales with amount}{6.6}
\practice{5}

A thermochemical equation carries its $\Delta H$ for the amounts written:
\[ \ce{CH4(g) + 2O2(g) -> CO2(g) + 2H2O(l)} \qquad
   \Delta H = -890.5~\text{kJ} \]
means \SI{890.5}{\kilo\joule} released per \emph{one mole} of \ce{CH4}.

\begin{itemize}[leftmargin=*,itemsep=0.3em]
  \item Double the coefficients $\Rightarrow$ $\Delta H$
        \answerblank[2cm]{doubles}.
  \item Reverse the reaction $\Rightarrow$ $\Delta H$
        \answerblank[3cm]{changes sign}.
  \item Half a mole of \ce{CH4} releases
        \answerblank[3cm]{\SI{445.3}{\kilo\joule}}.
\end{itemize}

\segment{INSTRUCTION B}{20}{Hess's law}

\notesection{Enthalpy is a state function}{6.9}
\practice{5}

\begin{center}
\fbox{\parbox{0.88\linewidth}{\centering
$\Delta H$ depends only on the initial and final states, \\
never on the route taken. \\[0.3em]
So reactions can be added, and their $\Delta H$ values add with them.}}
\end{center}

Two rules when manipulating:

\begin{enumerate}[leftmargin=*,itemsep=0.3em]
  \item Reverse a reaction $\Rightarrow$
        \answerblank[3.8cm]{flip the sign of $\Delta H$}.
  \item Multiply a reaction by $n$ $\Rightarrow$
        \answerblank[3.7cm]{multiply $\Delta H$ by $n$}.
\end{enumerate}

\begin{workedexample}[Worked example 5: finding an unmeasurable $\Delta H$]
Burning carbon straight to \ce{CO} cannot be done cleanly --- some
\ce{CO2} always forms. But these two \emph{can} be measured:

\begin{align*}
  (1)&\quad \ce{C(gr) + O2(g) -> CO2(g)}
     &\Delta H &= -393.5~\text{kJ} \\
  (2)&\quad \ce{CO(g) + \tfrac{1}{2}O2(g) -> CO2(g)}
     &\Delta H &= -283.0~\text{kJ}
\end{align*}

\textbf{Target:} $\ce{C(gr) + \tfrac{1}{2}O2(g) -> CO(g)}$.

Take (1) as written, and \emph{reverse} (2) so \ce{CO} ends up as a
product --- flipping its sign:
\begin{align*}
  \ce{C(gr) + O2 &-> CO2} &\Delta H &= -393.5 \\
  \ce{CO2 &-> CO + \tfrac{1}{2}O2} &\Delta H &= +283.0 \\
  \hline
  \ce{C(gr) + \tfrac{1}{2}O2 &-> CO} &\Delta H &= \mathbf{-110.5}~\text{kJ}
\end{align*}

\ce{CO2} cancels, and half the oxygen cancels. The answer matches the
tabulated $\Delta H_f^\circ$ of \ce{CO} exactly --- a good check.
\end{workedexample}

\segment{APPLICATION}{20}{Scaling and combining}

\begin{enumerate}[leftmargin=*,itemsep=0.4em]
  \item For $\ce{2CO(g) + O2(g) -> 2CO2(g)}$,
        $\Delta H = -566$~kJ. How much energy is released by
        \SI{1.00}{\mole} of \ce{CO}? \answerblank[3cm]{\SI{283}{\kilo\joule}}
  \item Give $\Delta H$ for $\ce{2CO2(g) -> 2CO(g) + O2(g)}$.
        \answerblank[3cm]{$+566$ kJ}
  \item A student combining reactions forgets to flip the sign when
        reversing one. If the correct answer is $-110.5$~kJ, what do they
        report, and what is the size of the error?
        \answerlines[2]{They report $-393.5 + (-283.0) = -676.5$~kJ. The
        error is $566$~kJ --- twice the magnitude of the reaction they
        mis-signed.}
\end{enumerate}

\begin{exitticket}
Why does Hess's law work --- what property of enthalpy makes it valid?
\keyonly{\begin{apcanswer}Enthalpy is a state function: its value depends
only on the initial and final states of the system, not on the path taken
between them. So any set of steps beginning and ending at the same places
must sum to the same $\Delta H$.\end{apcanswer}}
\end{exitticket}

% =====================================================================
\blocklesson{Bond and Formation Enthalpies}{CED 6.7, 6.8 \textbullet\ Zumdahl \S8.8, \S6.4}

\begin{objectives}
  \item Estimate $\Delta H$ from average bond enthalpies.
  \item Calculate $\Delta H^\circ$ from enthalpies of formation.
\end{objectives}

\reading{Zumdahl \S8.8}{412--415} \reading{Zumdahl \S6.4}{301--308}

\begin{warmup}[3]
  \item Reversing a reaction does what to $\Delta H$?
        \answerblank[3cm]{flips the sign}
  \item Breaking bonds is endothermic or exothermic?
        \answerblank[2.6cm]{endothermic}
  \item $\Delta H_f^\circ$ of \ce{O2(g)}: \answerblank[1.6cm]{0}
\end{warmup}

\segment{INSTRUCTION A}{25}{Two routes with opposite conventions}

\notesection{Bond enthalpies}{6.7}
\practice{5}

Breaking a bond always \answerblank[2.4cm]{requires} energy; forming one
always \answerblank[2.4cm]{releases} it. So:

\[ \boxed{\;\Delta H = \sum D(\text{bonds broken})
   - \sum D(\text{bonds formed})\;} \]

\begin{apctrap}
\textbf{Broken minus formed} --- the reverse of every other convention in
this unit, where it is products minus reactants. Bonds broken are on the
reactant side, so the order looks inverted, and students who have just
finished Hess's law reliably flip it.

A sanity check that costs nothing: if more energy is released forming bonds
than was spent breaking them, the reaction is exothermic and your answer
must come out \emph{negative}.
\end{apctrap}

\begin{workedexample}[Worked example 6: hydrogen and chlorine]
$\ce{H2(g) + Cl2(g) -> 2HCl(g)}$, using
$D(\ce{H-H}) = 432$, $D(\ce{Cl-Cl}) = 239$, $D(\ce{H-Cl}) = 427$
\si{\kilo\joule\per\mole}.

\textbf{Broken:} $432 + 239 = 671$~kJ. \\
\textbf{Formed:} $2(427) = 854$~kJ.
\[ \Delta H = 671 - 854 = \mathbf{-183}~\text{kJ} \]
Exothermic, as the negative sign should confirm.
\end{workedexample}

\notesection{Enthalpies of formation}{6.8}

$\Delta H_f^\circ$ is the enthalpy change forming \textbf{one mole} of a
compound from its elements in their standard states. By definition, an
element in its standard state has $\Delta H_f^\circ =$
\answerblank[1.4cm]{0}.

\[ \Delta H^\circ_{\text{rxn}} = \sum n\Delta H_f^\circ(\text{products})
   - \sum n\Delta H_f^\circ(\text{reactants}) \]

\begin{workedexample}[Worked example 7: the same reaction, two routes]
$\ce{CH4(g) + 2O2(g) -> CO2(g) + 2H2O(g)}$

\textbf{From $\Delta H_f^\circ$} ($\ce{CH4}$ $-75$, $\ce{CO2}$ $-393.5$,
$\ce{H2O(g)}$ $-242$, $\ce{O2}$ $0$):
\[ \Delta H^\circ = [-393.5 + 2(-242)] - [-75 + 0] =
   -877.5 + 75 = \mathbf{-802.5}~\text{kJ} \]

\textbf{From bond enthalpies:} broken $4(413) + 2(495) = 2642$;
formed $2(799) + 4(467) = 3466$; $\Delta H = \mathbf{-824}$~kJ.

\textbf{They disagree by about 22 kJ --- and that is expected.} Bond
enthalpies are \emph{averages} over many molecules and ignore the specific
molecular environment. Formation enthalpies are measured for the actual
substances. When a question gives you both, use $\Delta H_f^\circ$; use
bond enthalpies when that is all you have, and call the result an estimate.
\end{workedexample}

\segment{APPLICATION}{20}{Both routes}

\begin{enumerate}[leftmargin=*,itemsep=0.4em]
  \item Calculate $\Delta H^\circ$ for $\ce{N2 + 3H2 -> 2NH3}$ from
        $\Delta H_f^\circ(\ce{NH3}) = -46$~kJ/mol. \workspace[1.8cm]
        \keyonly{\begin{apcanswer}$2(-46) - 0 =
        \mathbf{-92}~\text{kJ}$\end{apcanswer}}
  \item Now estimate it from bond enthalpies
        ($\ce{N#N}$ 941, $\ce{H-H}$ 432, $\ce{N-H}$ 391).
        \workspace[2.4cm]
        \keyonly{\begin{apcanswer}Broken $941 + 3(432) = 2237$;
        formed $6(391) = 2346$;
        $\Delta H = \mathbf{-109}~\text{kJ}$ --- the same sign and rough
        size as $-92$, but not equal\end{apcanswer}}
  \item Account for the difference between your two answers.
        \answerlines[3]{Bond enthalpies are averaged over many different
        compounds, so they do not capture the actual N--H bonds in ammonia
        specifically. Formation enthalpies come from measurements on the
        real substances, so the $-92$~kJ value is the reliable one and the
        bond-enthalpy figure is an estimate.}
\end{enumerate}

\begin{exitticket}
State the two enthalpy expressions used in this block and say which one
runs backwards relative to the other.
\keyonly{\begin{apcanswer}Formation: products minus reactants. Bond
enthalpies: bonds \emph{broken} minus bonds \emph{formed} --- which is the
reactant side minus the product side, the reverse order. The reason is that
breaking bonds costs energy, so those terms must enter
positively.\end{apcanswer}}
\end{exitticket}

\end{document}
